\documentclass[11pt]{article}

\title{Technical description: Dual-decoder variational graph
autoencoder}
\author{}
\date{}

\usepackage{amsmath, amssymb}
\usepackage{float}
\usepackage[margin=2.5cm]{geometry}



\begin{document}
%% *******************************************************************

\section{Overview}
%% ===================================================================
%% -------------------------------------------------------------------
The model described here is a Dual-Decoder Variational Graph
Autoencoder (DualVGAE), trained jointly across a collection of
graphs derived from tabular data. Each tabular dataset is
converted into a graph in which rows become nodes and edges are
constructed by k-nearest-neighbour (kNN) similarity in feature
space. The model learns a low-dimensional latent representation
of each node that captures both local graph structure and
node-level feature content, and is trained to reconstruct both
the graph adjacency matrix and the original node features from
that representation.

The primary intended applications are unsupervised node
clustering and graph generation.


\section{Input Data and Graph Construction}
%% ===================================================================
%% -------------------------------------------------------------------

\subsection{Tabular preprocessing}
%% -------------------------------------------------------------------

Each source file (CSV, Parquet, NumPy array, or NPZ archive) is loaded
into a pandas DataFrame and processed as follows:
\begin{itemize}
\item Drop: non-informative columns (configurable): identifier
      columns, timestamps, and any columns designated as leakage.
\item One-hot encode categorical columns
      (configurable). High-cardinality categoricals should be handled
      upstream before this step.
\item Median imputation: of missing values, applied column-wise.
\item Standard scaling (zero mean, unit variance) using
      scikit-learn's StandardScaler, fit independently per graph. All
      downstream distances and GCN operations are performed on scaled
      features.
\end{itemize}
      

The result is a dense matrix $X$ of shape $[N, F]$, where $N$ is the
number of rows (nodes) and $F$ is the number of features after
encoding.

\subsection{Graph construction by kNN}
%% -------------------------------------------------------------------

An undirected graph $G = (V, E)$ is constructed from $X$ by connecting
each node to its k most similar neighbours in feature space
(default $k = 5$). Similarity is measured by cosine distance using
FAISS (IndexFlatIP on L2-normalised vectors), with a fallback to
scikit-learn's kneighbors\_graph when FAISS is unavailable.
      
The resulting directed kNN graph is symmetrised by taking the union of
forward and reverse edges $(A \gets A + A^\top)$, then deduplicated to produce a
simple undirected graph. The graph is stored as a PyTorch Geometric
Data object with node feature matrix x and adjacency encoded as a
sparse edge\_index tensor of shape [2, E].

No edge weights are used. The graph is treated as unweighted and
undirected throughout.



\section{Model Architecture}
%% ===================================================================
%% -------------------------------------------------------------------

The DualVGAE class extends PyTorch Geometric's VGAE base class with a
second decoder head for feature reconstruction. The full model
consists of three components: a shared GCN encoder, an
inner-product adjacency decoder (inherited from VGAE), and a new
MLP feature decoder.

\subsection{Encoder: VGAEEncoder}
%% -------------------------------------------------------------------
The encoder is a two-layer graph convolutional network implementing
the Kipf \& Welling (2017) spectral graph convolution (GCNConv). It
maps the pair (X, edge\_index) to a per-node posterior distribution
over a latent space of dimension $d$.

\subsubsection{Layer 1--Shared representation (conv\_shared)}
%% -------------------------------------------------------------------
A single GCNConv layer maps node features from F dimensions to H
dimensions (default H = 64), followed by a ReLU nonlinearity:
\[ h_i  =  \rho( W_{\text{shared}} \tilde{D}^{-1/2} \tilde{A}
\tilde{D}^{-1/2} X )_i, \]
where $\tilde{A} = A + I$ is the adjacency matrix with added
self-loops, $\tilde{D}$ is its degree matrix, 
$W_{\text{shared}} \in \mathbb{R}^{F×H}$
is a learned weight matrix, and $\rho(\cdot)$ is an activation
function (e.g. ReLU).
For each node $i$, this operation aggregates
the (degree-normalised) features of $i$ and all its immediate
neighbours, producing a neighbourhood-aware hidden representation
$h_i \in \mathbb{R}^H$. The receptive field at this layer is one hop.


\subsection{Layer 2--Parallel variational heads (conv\_mu,
conv\_logstd)}
%% -------------------------------------------------------------------

The hidden representation H is fed into two independent GCNConv layers
that share no weights, each mapping $H \to d$ (default $d = 32$). This
creates a Y-shaped split:
\begin{itemize}
\item conv\_mu: produces $\mu \in \mathbb{R}^{N \times d}$, the mean
of the posterior distribution $q(Z | X, A)$ for each node.

\item conv\_logstd: produces $\log \sigma \in \mathbb{R}^{N \times d}$,
the log standard deviation. Parameterising $\log \sigma$ rather than
$\sigma$ directly ensures positivity without constrained optimisation,
since $\sigma = \exp(\log \sigma)$ is always positive.
\end{itemize}

These two layers each perform a second round of neighbourhood
aggregation, giving the encoder an effective receptive field of two
hops: each node's latent distribution is informed by nodes up to
distance two in the graph.


\subsection{Reparameterisation}
%% -------------------------------------------------------------------

A latent embedding $z_i$ is sampled for each node using the
reparameterisation trick:
\[ z_i = \mu_i + \sigma_i \epsilon, \qquad \epsilon \sim N(0, I). \]
This allows gradients to flow through the sampling operation during
backpropagation: the randomness is isolated in $\epsilon$, which has
no learned parameters. The result is a matrix $Z \in \mathbb{R}^{N
\times d}$ of per-node latent embeddings.

At inference time (get\_embeddings), only $\mu$ is returned---the
deterministic mean embedding---rather than a stochastic sample,
which is standard practice for downstream tasks such as
clustering.

\subsubsection{Parameter count}
%% -------------------------------------------------------------------

\begin{table}[H]
\centering
\begin{tabular}{ r p{0.5\textwidth} }
conv\_shared & $W: [F \times H] + \text{bias}$: $[H] \to F H + H$ parameters \\
conv\_mu & $W: [H \times d] + \text{bias}$: $[d] \to H d + d$ parameters \\
conv\_logstd & $W: [H \times d] + \text{bias}$: $[d] \to H d + d$ parameters \\
Total (defaults: $F=10, H=64, d=32$)
& $10 \times 64 + 64 + 64 \times 32 + 32 + 64 \times 32 + 32 = 5,024$ parameters
\end{tabular}
\end{table}


\subsection{Adjacency decoder (inner product)}
%% -------------------------------------------------------------------
      
The adjacency decoder is the standard inner-product decoder inherited
from PyTorch Geometric's VGAE:
\[ \hat{A} = \phi( Z Z^\top ), \]
where $\phi$ is the sigmoid function and $\hat{A}_{ij} \in (0,1)$ is
the predicted probability of an edge between nodes $i$ and $j$. The
decoder has no learned parameters; all expressivity is in the
encoder. The predicted probability is compared against the
ground-truth adjacency using binary cross-entropy loss.

Because computing the full $N \times N$ matrix is infeasible for
graphs with hundreds of thousands of nodes, the loss is
evaluated only on a sampled set of positive edges (true edges in
the training split) and an equal number of negative samples
(randomly selected non-edges), following the standard VGAE
training protocol. 

\subsection{Feature decoder (MLP)}
%% -------------------------------------------------------------------

The feature decoder is a two-layer MLP that maps the latent embedding
$Z$ back to the input feature space, producing a reconstruction
$\hat{x}$:
\[ \hat{x} = W_2 \rho(W_1 Z + b_1) + b_2, \]
where $W_1 \in \mathbb{R}^{d \times H}$, $b_1 \in \mathbb{R}^H$,
$W_2 \in \mathbb{R}^{H \times F}$, and $b_2 \in \mathbb{R}^F$ are
learned parameters. The hidden dimension of the MLP is set to H (same
as the encoder's hidden dimension, default 64). The output has the
same shape as the input feature matrix X, enabling direct computation
of reconstruction error.

This decoder is an addition to the standard VGAE architecture and is
the key structural difference in DualVGAE. It provides a direct signal
to the encoder that node-level feature content should be preserved in
$Z$, not only graph topology.


\subsection{Training objective}
%% -------------------------------------------------------------------
The model is trained to minimise a weighted sum of three terms:
\[
L = \alpha L_{\text{edge}} + \beta L_{\text{feat}} + \gamma
L_{\text{KL}}
\]


\subsection{Edge reconstruction loss ($L_{\text{edge}}$)}
%% -------------------------------------------------------------------
Binary cross-entropy between the predicted adjacency probabilities
$\hat{A}$ and the ground-truth edge labels, evaluated on the
training edge split (positive edges) and an equal number of
randomly sampled negative edges:
\[
L_{\text{edge}} = -\frac{1}{\lvert E_{\text{train}} \rvert}
\sum_{i, j \in E_{\text{train}}} A_{ij} \log \hat{A}_{ij} +
(1 - A_{ij}) \log\big( 1 - \hat{A}_{ij} \big)
\]
This is the standard VGAE reconstruction loss. Minimising it
encourages the encoder to place nodes that are connected in the graph
near each other in latent space.


\subsection{Feature reconstruction loss ($L_{\text{feat}}$)}
%% -------------------------------------------------------------------
Mean squared error between the decoder's output $\hat{x}$ and the
original (scaled) node features $X$:
\[
L_{\text{feat}} = \frac{1}{N F} \sum_i \sum_f
\big( x_{if} - \hat{x}_{if} \big)^2.
\]
Minimising this loss encourages the encoder to retain sufficient
information in $Z$ to recover the node's own feature values, not just
its position in the graph. Because all features are standardised to
unit variance before training, the MSE is expressed in units of
standard deviations and is comparable across features.

\subsection{KL divergence loss ($L_{\text{KL}}$)}
%% -------------------------------------------------------------------
The Kullback-Leibler divergence between the learned posterior
$q(Z \mid X, A)$ and the standard normal prior $p(Z) = N(0, I)$,
normalised per node:
\begin{align*}
L_{\text{KL}}
&= \frac{1}{N} \mathcal{D}\big\{ q(Z \mid X, A) ~\Vert~ N(0, I) \big\}
\\
&= \frac{1}{2 N} \sum_i \sum_j ( \mu_{ij}^2 + \sigma_{ij}^2 - 2 \log
\sigma_{ij} - 1 )
\end{align*}
This regularisation term penalises the encoder for producing
posteriors that deviate far from the prior, which organises the
latent space into a smooth, continuous manifold and makes
generative sampling ($z \sim N(0, I)$) meaningful. The 1/N
normalisation is essential when training on graphs of varying
sizes, as it prevents large graphs from contributing
disproportionately to the KL term.

\subsection{Loss weights}
%% -------------------------------------------------------------------
The weights $\alpha, \beta, \gamma$ are user-configurable
hyperparameters (default 1.0 each). Their primary purpose is to
balance the three loss terms when they differ in magnitude at
initialisation. A recommended starting procedure is to inspect the
per-term losses logged at the first epoch and scale the weights so
that all three terms are roughly equal before tuning for downstream
task performance. Increasing $\alpha$ relative to $\beta$ emphasises
graph structure preservation; increasing $\beta$ emphasises feature
fidelity.


\section{Optimisation and training protocol}
%% ===================================================================
%% -------------------------------------------------------------------

\subsection{Optimizer}
%% -------------------------------------------------------------------
All parameters (encoder GCN layers and feature decoder MLP) are
trained jointly using the Adam optimiser with a fixed learning
rate (default $1 \times 10^{-3}$) and PyTorch default
$\beta_1 = 0.9, \beta_2 = 0.999$. No learning rate schedule or weight
decay is applied by default. 

\subsection{Edge split}
%% -------------------------------------------------------------------
For each graph, PyTorch Geometric's RandomLinkSplit transform holds
out a fraction of edges as a validation set (default 5\%) and a test
set (default 5\%) before training begins. The model is trained on the
remaining edges only. Held-out edges are used to compute AUC and
Average Precision at validation epochs, providing an unbiased estimate
of edge reconstruction quality.

\subsection{Multi-graph training}
%% -------------------------------------------------------------------
The model is trained sequentially across all graphs in each epoch,
rather than batching nodes across graphs. This is necessary
because the GCN convolution is defined with respect to each
graph's own adjacency structure, which cannot be trivially mixed
across graphs. Within a single graph, PyTorch Geometric's
DataLoader is used to form mini-batches; graphs exceeding a
configurable node threshold (default $300,000$ nodes) are
processed one at a time (batch size $1$) to limit GPU memory
consumption.
      
All graphs are loaded from disk on first encounter and cached in RAM
for subsequent epochs. Given the available hardware
configuration ($\approx 240$ GB RAM), all graphs in the dataset are
expected to fit simultaneously in memory after the first epoch,
eliminating repeated disk I/O.

\subsection{Gradient checkpointing}
%% -------------------------------------------------------------------
When training on a CUDA device, gradient checkpointing is applied to
the shared GCN layer (conv\_shared). Rather than storing
intermediate activations in GPU memory for use during
backpropagation, the forward pass of conv\_shared is recomputed
during the backward pass. This trades approximately 30\%
additional compute for approximately 40\% reduction in GPU
activation memory, and is particularly beneficial for the
largest graphs.


\section{Model outputs}
%% ===================================================================
%% -------------------------------------------------------------------

\subsection{Node embeddings}
%% -------------------------------------------------------------------
After training, the encoder is applied to each graph in inference mode
(no sampling; $Z = \mu$) to produce a matrix of node embeddings
of shape $[N, d]$. Embeddings are exported per graph as CSV files
(one row per node, $d$ columns $z_0$ through $z_{d-1}$). These
embeddings are the primary output for downstream tasks.

\subsection{Generated graphs}
%% -------------------------------------------------------------------
New graphs can be generated by sampling latent vectors
$z \sim N(0, I)$ and passing them through the inner-product
decoder. The result is a probabilistic adjacency matrix $\hat{A}$;
thresholding at a chosen probability (default 0.5) gives a binary
adjacency matrix. Generated graphs have no associated node features
unless a separate feature generation step is added.


\subsection{Training loss history}
%% -------------------------------------------------------------------
All four loss values (total, edge, feature, KL) are recorded per epoch
and written to a CSV file at the end of training. This breakdown
facilitates post-hoc diagnosis of imbalanced loss terms and
tuning of $\alpha, \beta, \gamma$. 



\section{Evaluation}
%% ===================================================================
%% -------------------------------------------------------------------

\subsection{Adjacency reconstruction}
%% -------------------------------------------------------------------
Edge reconstruction quality is measured using two metrics computed on
the held-out validation edges and an equal number of randomly sampled
negative edges:

\begin{itemize}
\item AUC (Area Under the ROC Curve): measures the model's ability to
  rank true edges above non-edges. A random model scores $0.5$; a
  well-trained VGAE typically achieves $0.85--0.95$ on structured
  graph data.
\item Average Precision (AP): the area under the
  precision-recall curve, which is more informative than AUC when
  the graph is sparse (far more non-edges than edges), as is
  typical of kNN-constructed graphs.
\end{itemize}


\subsection{Feature reconstruction}
%% -------------------------------------------------------------------
For each graph and each feature dimension, the following metrics are
computed between the original scaled features X and the decoder
output $\hat{x}$:

\begin{table}[H]
\centering
\begin{tabular}{ r c p{0.5\textwidth} }
Metric & Formula & Interpretation \\
\hline
MSE
& $\frac{1}{N} \lVert x - \hat{x} \rVert^2$
& Mean squared error in units of standard deviations (features are
standardised). Lower is better. \\
MAE
& $\frac{1}{N} \sum_i \lvert x_i - \hat{x}_i \rvert$
& Mean absolute error; less sensitive to outliers than MSE. \\
$R^2$
& $1 - \frac{\text{SS}_{\text{res}}}{\text{SS}_{\text{tot}}}$
& Fraction of variance explained by the decoder.
\end{tabular}
\end{table}

Results are aggregated across graphs and summarised per feature
dimension, enabling identification of which features are
well-preserved in the latent space and which are discarded. Low $R^2$
for a feature is not necessarily a failure: it may indicate that the
feature is idiosyncratic to individual nodes rather than correlated
with local graph structure, and is therefore not recoverable from a
neighbourhood-aware embedding.


\subsection{Embedding quality}
%% -------------------------------------------------------------------
Node embeddings are evaluated for clustering quality using:

\begin{itemize}
\item Silhouette score $\in [-1, 1]$: measures cohesion within
  clusters relative to separation between clusters. Values above
  $0.3$ indicate meaningful structure; above $0.5$ is strong. Computed
  on a random subsample of nodes ($\leq 10,000$) due to
  $\mathcal{O}(N^2)$ complexity.
\item Davies-Bouldin index: ratio of within-cluster scatter to
  between-cluster separation; lower is better. Computed on the
  full embedding as it scales as $\mathcal{O}(N k)$.
\item $k$ sweep: both metrics are evaluated across a range of $k$
  values for $k$-means clustering, with the optimal $k$ identified as
  the peak in silhouette score and the elbow in inertia.
\end{itemize}
  

Prior coverage of the latent space is assessed by verifying that the
mean and standard deviation of $Z$ across all nodes and all graphs are
close to 0 and 1 respectively, matching the $N(0, I)$ prior from which
generated graphs are sampled.



\section{Tuning parameters}
%% ===================================================================
%% -------------------------------------------------------------------

\begin{table}[H]
\begin{tabular}{ r p{0.65\textwidth} }
$k$
& Number of nearest neighbours per node in kNN graph
construction. Default: 5. Increase for denser graphs; decrease
if the model overfits to local structure. \\

hidden\_channels ($H$)
& Dimensionality of the shared GCN hidden layer and the MLP feature
decoder's hidden layer. Default: 64. \\

latent\_dim ($d$)
& Dimensionality of the latent space. Default: 32. Increase if
clustering results are indistinct; decreasing improves generation
regularity. \\

loss\_alpha ($\alpha$)
& Weight on the edge reconstruction BCE loss. Default: 1.0. \\

loss\_beta ($\beta$)
& Weight on the feature reconstruction MSE loss. Default: 1.0. \\

loss\_gamma ($\gamma$)
& Weight on the KL divergence term. Default: 1.0. Increasing $\gamma$
produces a more organised latent space at the cost of reconstruction
fidelity. \\

epochs
& Number of full passes through the graph collection. Default: 200. \\

lr
& Adam learning rate. Default: $1 \times 10^{-3}$. \\

batch\_size
& Number of graphs per DataLoader batch (for graphs below
large\_graph\_threshold). Default: 4. \\

large\_graph\_threshold
& Graphs with more nodes than this value are processed with
batch\_size=1 regardless of the batch\_size setting. Default:
300,000. \\

val\_fraction
& Fraction of edges held out for validation. Default: 0.05. \\

test\_fraction
& Fraction of edges held out for testing. Default: 0.05.
\end{tabular}
\end{table}


\section{Dependencies}
%% ===================================================================
%% -------------------------------------------------------------------

\begin{table}[H]
\begin{tabular}{ r p{0.65\textwidth} }
\texttt{PyTorch}
& Tensor operations, autograd, Adam optimiser \\
\texttt{PyTorch Geometric}
& GCNConv, VGAE base class, Data/DataLoader, RandomLinkSplit \\
\texttt{FAISS}
& Approximate kNN graph construction (CPU or GPU) \\
\texttt{scikit-learn}
& StandardScaler, KMeans, silhouette\_score, davies\_bouldin\_score; kneighbors\_graph fallback \\
\texttt{pandas/NumPy}
& Tabular data loading, preprocessing, CSV output \\
\texttt{psutil}
& Runtime memory availability reporting
\end{tabular}
\end{table}


%% *******************************************************************
\end{document}

